\documentclass[12pt,a4paper]{article}
\usepackage[utf8]{inputenc}
\usepackage[korean]{babel}
\usepackage{kotex}
\usepackage{geometry}
\usepackage{amsmath}
\usepackage{amsfonts}
\usepackage{amssymb}
\usepackage{graphicx}
\usepackage{fancyhdr}
\usepackage{setspace}
\usepackage{titlesec}
\usepackage{cite}
\usepackage{url}
\usepackage{hyperref}
\usepackage{xcolor}
\usepackage{listings}
\usepackage{enumitem}
\usepackage{multicol}
\usepackage{needspace}
\usepackage{afterpage}
\usepackage{fontawesome5}

% 페이지 설정
\geometry{margin=2.5cm}
\setstretch{1.5}

% 폰트 설정
\setmainhangulfont[BoldFont={D2Coding}]{D2Coding}
\setmonofont{JetBrains Mono}

% 헤더와 푸터 설정
\pagestyle{fancy}
\fancyhf{}
\fancyhead[L]{김찬규 | Portfolio}
\fancyhead[R]{\footnotesize\icon{\faEnvelope}{10pt}{\href{mailto:kjkj5868@gmail.com}{kjkj5868@gmail.com}} | \icon{\faGithub}{10pt}{\href{https://github.com/chan99k}{github.com/chan99k}}}
\fancyfoot[C]{\thepage}

% 제목 스타일 설정 및 페이지 나누기 방지
\titleformat{\section}[block]{\Large\bfseries}{\thesection}{1em}{}
\titlespacing*{\section}{0pt}{20pt plus 2pt minus 2pt}{12pt plus 2pt minus 2pt}
\titleformat{\subsection}{\large\bfseries}{\thesubsection}{1em}{}
\titleformat{\subsubsection}{\normalsize\bfseries}{\thesubsubsection}{1em}{}

% 섹션 넘버링 깊이 설정
\setcounter{secnumdepth}{3}
\setcounter{tocdepth}{3}

% 2단 환경에서의 제목 스타일 (더 작은 폰트)
\newcommand{\multicolsubsection}[1]{%
    \stepcounter{subsubsection}%
    \addcontentsline{toc}{subsubsection}{\protect\numberline{\thesubsubsection}#1}%
    \vspace{0.5em}%
    {\small\bfseries \thesubsubsection.\ #1}%
    \vspace{0.3em}%
}

\newcommand{\multicolsubsubsection}[1]{%
    \vspace{0.3em}%
    {\small\bfseries #1}%
    \vspace{0.2em}%
}

% 페이지 나누기 패널티 조정
\clubpenalty=10000
\widowpenalty=10000
\raggedbottom

% 코드 스타일 설정
\lstset{
    basicstyle=\ttfamily\small,
    backgroundcolor=\color{gray!10},
    frame=single,
    breaklines=true,
    showstringspaces=false
}

% 하이퍼링크 색상 설정
\hypersetup{
    colorlinks=true,
    linkcolor=black,
    filecolor=magenta,
    urlcolor=black,
}

% 아이콘 명령어 정의
\newcommand{\icon}[3]{\makebox[#2][c]{\textcolor{black}{#1}}\hspace{0.2em}#3}

% 구분선 명령어 정의
\newcommand{\divider}{%
    \vspace{0.5cm}%
    \begin{center}%
    \textcolor{lightgray}{\rule{0.2\textwidth}{0.5pt}}%
    \end{center}%
    \vspace{0.5cm}%
}

\newcommand{\thickdivider}{%
    \vspace{0.5cm}%
    \begin{center}%
    \textcolor{gray}{\rule{0.3\textwidth}{1pt}}%
    \end{center}%
% \vspace{0.1cm}%
}

\newcommand{\dotdivider}{%
    \vspace{0.5cm}%
    \begin{center}%
    \textcolor{lightgray}{$\bullet \quad \bullet \quad \bullet$}%
    \end{center}%
    \vspace{0.5cm}%
}

\newcommand{\stardivider}{%
    \vspace{0.5cm}%
    \begin{center}%
    \textcolor{lightgray}{$\star \quad \star \quad \star$}%
    \end{center}%
    \vspace{0.5cm}%
}

\begin{document}

% 제목 페이지
    \begin{titlepage}
        \centering
        \vspace*{2cm}

        {\huge\bfseries Portfolio\par}
        \vspace{0.5cm}
        {\Large 개인 프로젝트 경험과 문제 해결 과정\par}

        \vspace{2cm}
        {\large 김찬규\par}

        \vspace{1cm}
        {\large \today\par}

        \vfill

        {\large
        Java, Spring Boot, Spring Batch, MySQL 등을 활용한\\
        백엔드 시스템 설계와 개발 경험
        \par}

    \end{titlepage}

%----------------------------------------------------------------------------------------
%	목차
%----------------------------------------------------------------------------------------

    \tableofcontents
    \newpage

%----------------------------------------------------------------------------------------
%	초록 | 필요없으니 생략
%----------------------------------------------------------------------------------------

    % \begin{abstract}
    % [여기에 포트폴리오 요약 내용을 작성하세요. 주요 프로젝트와 기술적 성과를 간략히 기술합니다.]
    % \end{abstract}

%----------------------------------------------------------------------------------------
%	INTRODUCTION
%----------------------------------------------------------------------------------------
    \section{Introduction}

    [여기에 서론 내용을 작성하세요. 포트폴리오의 목적과 배경을 설명합니다.]

    \subsection{Personal Summary}

    \begin{minipage}[t]{\textwidth}
        \setlength{\baselineskip}{16pt} % 줄간격 조절
        \fontsize{10}{14}\selectfont    % 폰트 크기 10pt, 줄간격 14pt
        \hspace*{0.5em} \textbf{Java/Spring 기반의 백엔드 개발자}로서, 개인 프로젝트에서 \textbf{API 응답 시간과 같은 지표를 분석}하여 병목 지점을 찾아 \textbf{API 성능을 최적화}하고, \textbf{Spring Batch}를 활용해 \textbf{데이터를 처리하는 과정을 자동화}한 경험이 있습니다. 나아가, 이러한 개발 과정을 \textbf{Github Actions 기반의 CI 파이프라인}으로 자동화하여 구축해본 경험이 있습니다.
    \end{minipage}

    \thickdivider
%----------------------------------------------------------------------------------------
%	프로젝트 1
%----------------------------------------------------------------------------------------
    \section{WashFit | 개인 프로젝트}

    \subsection{프로젝트 배경 및 목표}

    [여기에 WashFit 프로젝트의 배경과 목표를 작성하세요.]

    \subsection{사용 기술}

    [여기에 사용된 기술 스택을 작성하세요.]

    \divider

    \subsection{문제 및 해결방법}

    % 문제 해결 섹션을 2단으로 시작
    \begin{multicols}{2}
        \setlength{\columnsep}{2.5cm}
        \setlength{\columnseprule}{0pt}

        \multicolsubsection{멀티 모듈 프로젝트 아키텍처 설계 및 구현}

        \multicolsubsubsection{문제 상황}

        [여기에 멀티 모듈 아키텍처 도입 전 문제 상황을 작성하세요.]

        \multicolsubsubsection{해결 방안}

        [여기에 멀티 모듈 아키텍처 해결 방안을 작성하세요.]

        \begin{figure}[htbp]
            % \centering
            % \includegraphics[width=0.6\textwidth]{images/multimodule-architecture}
            % \caption{멀티 모듈 아키텍처 구조도}
            % \label{fig:multimodule}
        \end{figure}

        \begin{lstlisting}[caption=멀티 모듈 구조]
"washfit-backend/
├── module-api/          # REST API 엔드포인트
├── module-batch/        # Spring Batch 처리
├── module-domain/       # 도메인 로직
├── module-common/       # 공통 유틸리티
└── module-admin/        # 관리자 기능
        "
        \end{lstlisting}

        \multicolsubsubsection{구현 결과}

        [여기에 멀티 모듈 아키텍처 구현 결과와 성과를 작성하세요.]

        \multicolsubsection{QueryDSL을 활용한 동적 쿼리 및 타입 안정성 확보}

        \multicolsubsubsection{문제 상황}

        [여기에 QueryDSL 도입 전 문제 상황을 작성하세요.]

        \multicolsubsubsection{해결 방안}

        [여기에 QueryDSL 해결 방안을 작성하세요.]

        \multicolsubsubsection{구현 결과}

        [여기에 QueryDSL 구현 결과와 성과를 작성하세요.]

        \multicolsubsection{Spring Batch 모듈 설계를 통한 데이터 처리 아키텍처 구축}

        \multicolsubsubsection{문제 상황}

        [여기에 Spring Batch 도입 전 문제 상황을 작성하세요.]

        \multicolsubsubsection{해결 방안}

        [여기에 Spring Batch 해결 방안을 작성하세요.]

        \multicolsubsubsection{구현 결과}

        [여기에 Spring Batch 구현 결과와 성과를 작성하세요.]

    \end{multicols}
    % WashFit 프로젝트 2단 종료

    \needspace{5\baselineskip}
    \section{KernelEngine | 개인 프로젝트}

    \subsection{프로젝트 배경 및 목표}

    KernelEngine은 기술 블로그 게시글을 주기적으로 수집하고 키워드 검색 기능을 제공하는 웹 서비스입니다. 2023년 11월부터 12월까지 3인 팀 프로젝트로 진행되었으며, 자동화된 데이터 수집과 효율적인 검색 시스템 구축에 중점을 두었습니다.

    \subsection{사용 기술}

    [여기에 사용된 기술 스택을 작성하세요.]


    \divider

    \subsection{문제 및 해결방법}

    % KernelEngine 문제 해결 섹션을 2단으로 시작
    \begin{multicols}{2}
        \setlength{\columnsep}{2.5cm}
        \setlength{\columnseprule}{0pt}

        \multicolsubsection{MySQL FULLTEXT를 활용한 검색 시스템 구현}

        \multicolsubsubsection{문제 상황}

        [여기에 검색 시스템 구현 전 문제 상황을 작성하세요.]

        \multicolsubsubsection{해결 방안}

        [여기에 MySQL FULLTEXT 검색 해결 방안을 작성하세요.]

        \multicolsubsection{RSS 파서를 활용한 블로그 크롤링 시스템 구현}

        \begin{figure}[htbp]
            \centering
            \includegraphics[width=0.8\textwidth]{images/rss-crawler-system}
            \caption{RSS 크롤링 시스템 아키텍처}
            \label{fig:rss-crawler}
        \end{figure}

        \multicolsubsubsection{문제 상황}

        [여기에 RSS 크롤링 시스템 구현 전 문제 상황을 작성하세요.]

        \multicolsubsubsection{해결 방안}

        [여기에 RSS 크롤링 시스템 해결 방안을 작성하세요.]

        \multicolsubsubsection{구현 결과}

        [여기에 RSS 크롤링 시스템 구현 결과와 성과를 작성하세요.]

        \multicolsubsection{Spring Batch를 활용한 데이터 처리 및 통계 관리 시스템}

        \multicolsubsubsection{문제 상황}

        [여기에 Spring Batch 데이터 처리 시스템 구현 전 문제 상황을 작성하세요.]

        \multicolsubsubsection{해결 방안}

        [여기에 Spring Batch 데이터 처리 시스템 해결 방안을 작성하세요.]

        \begin{figure}[htbp]
            \centering
            \includegraphics[width=0.7\textwidth]{images/spring-batch-architecture}
            \caption{Spring Batch 데이터 처리 파이프라인}
            \label{fig:spring-batch}
        \end{figure}

    \end{multicols}
    % KernelEngine 프로젝트 2단 종료

    \needspace{5\baselineskip}

    \section{Conclusion and Future Plans}

    \subsection{주요 성과}

    [여기에 포트폴리오의 주요 성과를 작성하세요.]

    \subsection{기술적 성장}

    [여기에 기술적 성장 과정에 대한 내용을 작성하세요.]

    \subsection{향후 학습 방향}

    [여기에 향후 학습 방향과 목표를 작성하세요.]

    \section*{}

    \begin{thebibliography}{9}

        \bibitem{WashFit}
        WashFit Github Repository.
        \textit{https://docs.spring.io/spring-boot/docs/current/reference/htmlsingle/}

        \bibitem{Kernel Engine}
        KernelEngine Github Repository.
        \textit{http://querydsl.com/static/querydsl/latest/reference/html/}


    \end{thebibliography}

\end{document}